\documentclass[11pt]{amsart}
\usepackage{amsmath,amssymb,amsthm}
\usepackage{graphicx}
\usepackage{hyperref}

\newtheorem{theorem}{Theorem}
\newtheorem{corollary}{Corollary}
\theoremstyle{definition}
\newtheorem{definition}{Definition}

\DeclareMathOperator{\Aff}{Aff}
\setlength{\emergencystretch}{3em}

\title[Peak-point positive-isometry extension]{A peak-point partial answer for order-preserving isometry extensions}
\author{Agent lane 16}
\date{2026-06-29}

\begin{document}
\maketitle

\section*{Status}
This is a candidate partial result for the final Problem in Tanmoy Paul and
T. S. S. R. K. Rao, \emph{Order-preserving unique Hahn-Banach extensions},
arXiv:2504.03386~\cite{PaulRao}.  The result addresses the literal reading in
which an order-preserving isometric embedding
\[
  \Phi:A(K)\longrightarrow C(K)
\]
is to have a unique positive linear isometric extension
\[
  T:C(K)\longrightarrow C(K).
\]
It does not settle the broader projection-valued or operator-space variants
suggested by the surrounding paragraph in the source paper.

\section*{Source problem}
The source problem appears on PDF page 14, in the final displayed Problem of
\cite{PaulRao}.  A crop of the statement is included below.

\begin{center}
\includegraphics[width=\linewidth]{figures/open_problem_crop.png}
\end{center}

In transcription, the final sentence asks:
\begin{quote}
In view of the operator versions from Rao, it would be interesting to find
conditions under which an order-preserving isometric embedding
\(\Phi:A(K)\to C(K)\) has unique extension to an into order-preserving isometry
from \(C(K)\) to \(C(K)\).
\end{quote}

\section*{Definitions}
Throughout, \(K\) is a compact convex subset of a locally convex real
topological vector space, \(A(K)\) denotes the real Banach space of continuous
affine real-valued functions on \(K\), and \(C(K)\) denotes the real Banach
lattice of continuous real-valued functions on \(K\), both with the supremum
norm.

\begin{definition}
A point \(x\in K\) is a peak point for \(C(K)\) if there is \(p_x\in C(K)\)
such that \(0\leq p_x\leq 1\), \(p_x(x)=1\), and \(p_x(y)<1\) for every
\(y\neq x\).
\end{definition}

Every compact metrizable Hausdorff space has this property: for example, if
\(d\) is a compatible metric, \(p_x(y)=1-\min\{d(x,y),1\}/2\) peaks at \(x\).

\section*{Proof Intuition}
Let \(T:C(K)\to C(K)\) be a positive unital extension of a composition
embedding \(a\mapsto a\circ h\).  Evaluating \(Tf\) at a point \(s\) gives, by
the Riesz-Markov representation theorem, integration against a probability
measure \(\mu_s\) on \(K\).  Because \(T\) extends \(a\mapsto a\circ h\), the
measure \(\mu_s\) represents the point \(h(s)\) on affine functions.  The
isometry condition then forces these representing measures to be point masses:
if a continuous peak function at \(h(s)\) has norm one, so does its image under
\(T\), and at a point where this norm is attained the associated probability
measure must be supported entirely at \(h(s)\).  The affine separation of
points and injectivity of \(h\) move that conclusion back to the original
point \(s\).  Thus each \(\mu_s=\delta_{h(s)}\), so \(T\) is necessarily the
composition operator \(f\mapsto f\circ h\).

The part of the source problem that remains open is the general classification
of embeddings \(\Phi:A(K)\to C(K)\) for which this uniqueness holds, especially
under interpretations closer to the projection or operator-space language in
the source paragraph.

\section*{The partial theorem}
\begin{theorem}
Let \(K\) be a compact convex set such that every point of \(K\) is a peak
point for \(C(K)\).  Let \(h:K\to K\) be a homeomorphism, and define
\[
  \Phi:A(K)\longrightarrow C(K),\qquad \Phi(a)=a\circ h.
\]
Then \(\Phi\) is a unital order-preserving linear isometry.  Moreover, it has a
unique positive linear isometric extension \(T:C(K)\to C(K)\), namely
\[
  T(f)=f\circ h\qquad (f\in C(K)).
\]
\end{theorem}

\begin{proof}
Since \(h\) is onto, the map \(a\mapsto a\circ h\) preserves suprema of
absolute values and hence is an isometry from \(A(K)\) into \(C(K)\).  It is
linear, unital, and order preserving.  The composition operator
\[
  C_h:C(K)\to C(K),\qquad C_h(f)=f\circ h,
\]
is a positive linear isometry extending \(\Phi\), so existence is clear.

It remains to prove uniqueness.  Let \(T:C(K)\to C(K)\) be any positive linear
isometry with \(T|_{A(K)}=\Phi\).  Since \(1\in A(K)\) and \(\Phi(1)=1\), we
have \(T1=1\).  For each \(s\in K\), the functional
\[
  f\longmapsto Tf(s)
\]
is a positive norm-one functional on \(C(K)\).  By the Riesz-Markov
representation theorem~\cite[Theorem 7.17]{Rudin}, there is a probability
measure \(\mu_s\) on \(K\) such that
\[
  Tf(s)=\int_K f\,d\mu_s\qquad(f\in C(K)).
\]
For every \(a\in A(K)\),
\[
  \int_K a\,d\mu_s = T a(s)=\Phi(a)(s)=a(h(s)).
\]
Thus \(\mu_s\) represents the point \(h(s)\) on affine functions.

Fix \(s\in K\), and choose a peak function \(p\in C(K)\) for the point
\(h(s)\), with \(0\leq p\leq 1\), \(p(h(s))=1\), and \(p<1\) off
\(\{h(s)\}\).  Since \(T\) is an isometry and \(\|p\|=1\), the positive
continuous function \(Tp\) has norm \(1\).  Compactness gives a point \(r\in K\)
with
\[
  1=Tp(r)=\int_K p\,d\mu_r.
\]
Because \(0\leq p\leq 1\), and \(p=1\) only at \(h(s)\), this forces
\(\mu_r=\delta_{h(s)}\).

On the other hand, \(\mu_r\) represents \(h(r)\) on affine functions, so for
every \(a\in A(K)\),
\[
  a(h(r))=\int_K a\,d\mu_r=a(h(s)).
\]
Continuous affine functions separate points of \(K\), hence \(h(r)=h(s)\).
Since \(h\) is injective, \(r=s\).  Therefore \(\mu_s=\delta_{h(s)}\).

As \(s\in K\) was arbitrary, for every \(f\in C(K)\) and \(s\in K\),
\[
  Tf(s)=\int_K f\,d\mu_s=f(h(s)).
\]
Thus \(T=C_h\), proving uniqueness.
\end{proof}

\begin{corollary}
If \(K\) is compact metrizable and convex, then for every homeomorphism
\(h:K\to K\), the embedding \(a\mapsto a\circ h\) from \(A(K)\) into \(C(K)\)
has the unique positive linear isometric extension \(f\mapsto f\circ h\).
In particular, for \(h=\mathrm{id}_K\), the canonical inclusion
\(A(K)\subseteq C(K)\) has the identity as its unique positive linear
isometric extension under this literal \(C(K)\to C(K)\) reading.
\end{corollary}

\section*{Verification and novelty notes}
The proof is purely analytic; no finite search or numerical verification is
used.  The only external theorem invoked is the standard Riesz-Markov
representation theorem for positive functionals on \(C(K)\).

Before promotion, the run's cheap indexes were searched for arXiv:2504.03386,
the paper title, and core phrases around order-preserving Hahn-Banach and
isometry extensions; no existing packet was found.  A bounded web search for exact
phrases from the source problem and for Rao's operator-version wording found
no direct later answer.  Novelty confidence is therefore moderate for the
scoped partial theorem, with the main human-review risk being interpretation
of the source problem rather than the proof itself.

\section*{Human-review recommendation}
Review as a scoped partial answer.  The proof should be checked for the
measure-collapse step and for the stated point-peak hypothesis.  The result
should not be treated as a full answer to the source problem unless a reviewer
accepts the restricted composition-embedding setting as the intended class.

\begin{thebibliography}{9}
\bibitem{PaulRao}
Tanmoy Paul and T. S. S. R. K. Rao,
\emph{Order-preserving unique Hahn-Banach extensions},
arXiv:2504.03386, 2025.
\url{https://arxiv.org/abs/2504.03386}

\bibitem{Rudin}
Walter Rudin,
\emph{Real and Complex Analysis},
3rd edition, McGraw-Hill, 1987.
\end{thebibliography}

\end{document}
